\chapter{语义单元：有限分辨率、形质绑定与身份保持}
\label{chp:quantization_reality}

前一篇建立了目标、结构与输入共同作用的状态动力学。现在我们需要确定这些更新究竟作用于什么：一个词表索引、一个持续跟踪的对象、一组分布式活动，还是它们之间的接口？如果不区分这些对象，状态变化便容易被误读为现实对象的变化，内部表示的连续性也容易被当成事实的正确性。我们在本章把语义单元定义为具有身份、关系位置、属性内容与证据来源的任务相关记录，并研究记录如何从观测中形成、如何绑定、如何在表示变换中保持可检查的连续性。

我们使用有限分辨率编码解释离散化，而不把认知计算建立在量子统计的类比上。形与质分别描述关系组织和属性内容，二者的联合既可以由显式记录实现，也可以由张量或分布式编码实现；哪一种实现有效，需要通过任务、资源与错误代价判断。本章先规定单元的类型和绑定条件，再讨论传播、近似不变量、属性谱系及跨系统比较。由此得到的是一组可实现、可反驳的计算约束，不是关于所有智能载体的唯一微观实体假说。

\section{从输入索引到具有身份的联合记录}
\label{sec:section_3160}

我们首先区分原始索引和语义单元。文本中的词表编号用于定位某段输入编码，视觉块编号用于定位一块采样区域，数据库键用于定位某条记录。它们都可以称为 Token，但索引的存在本身不说明系统已经知道被索引的对象是什么。同一个词表编号可能出现在两个不同人物的发言中，同一个图像块可能同时包含桌面和杯沿，不同数据库键也可能指向同一现实对象。语义处理的第一项责任因此是声明索引与对象之间的关系，而不是赋予每个索引一种独立实体地位。

我们把一个语义单元写成带类型的记录。设身份集合为有限或可数的集合，关系上下文为一张带类型图，属性空间为按任务规定的有限维空间，证据集合包含来源、采样时刻和处理版本。身份是系统用于持续指称对象的键，不是对现实对象绝对同一性的证明。图节点可以表示对象、事件或角色，但类型必须显式区分。属性可以是颜色、文本向量、测量数值、动作可行性或情绪效价，异类字段不共享同一种单位。缺失值需要专门的标记，不能用零数值兼任未知，因为零温差和未知温差会触发不同的判断。

在固定任务与固定结构版本下，我们定义记录为
\[
 \mathcal S_i=(\iota_i,\mu_i,q_i,a_i,\mathcal P_i),\qquad
 \mu_i\in\mathbb R^{d_s},\quad q_i\in\mathbb R^{d_q},\quad a_i\in[0,1].
\]
其中身份键为 $\iota_i$，结构编码为 $\mu_i$，属性编码为 $q_i$，活动门控为 $a_i$，来源记录为 $\mathcal P_i$。两个向量在进入联合模型前按规定尺度归一化为无量纲数；实际测量值及其单位另存于有类型字段中。活动门控表示当前参与计算的程度，不是对象存在概率，更不是以焦耳计量的能耗。我们保留 Semantion 一词作为这种语义单元的名称，形元和质元则指结构分量与属性分量；这些名称不附带任何基本粒子的统计假设。

若任务需要结构与内容的双线性交互，我们可以定义联合矩阵 $B_i=\mu_iq_i^{\mathsf T}$，它属于 $d_s\times d_q$ 的实矩阵空间；向量化后等价于选定基底下的张量积表示。这里张量积是编码定义，不是所有绑定的必然形式。对非零标量，把结构向量乘以该标量并把属性向量除以同一标量，联合矩阵保持不变，因此单凭联合矩阵不能恢复两个分量的唯一尺度。若某一分量为零，整个矩阵也为零，属性来源和身份便不能靠它保留。显式记录、归一化条件或额外读出接口因此仍然必要，不能因矩阵维度较高就声称保存了对象的一切性质。

例如，在桌面任务中，系统看到一个红色杯子。我们把杯子的身份、支撑关系和抓取角色放入关系记录，把颜色、边缘形状及测量置信放入属性记录。图像中的颜色片段不自动成为杯子的属性，必须先通过观测区域、时间窗口和对象关联建立归属。杯子移动后，空间坐标改变而身份可以继续；杯子破损后，一些属性改变而来源历史仍然保留；杯子被另一个外观相似的杯子替换时，视觉相似不能单独证明身份延续。这些情况说明，最小计算单元的粒度取决于任务需要追踪的差异，不能事先规定为一个词、一个像素或一个固定长度向量。

局部结构网络和全局分布表示也在此处分工。前者显式组织对象、部件、邻接与事件角色，适合检查关系约束；后者把跨对象、跨情境的统计信息分布在共享参数和活动向量中，适合相似检索与泛化。我们通过身份键、版本、绑定映射和读出函数连接二者，而不要求它们互相替代。局部图可以为一个对象提供多个关系位置，全局向量也可以同时携带多个候选含义，二者的对应需要由任务接口确认。脑区功能的相似描述可以启发接口设计，但不能证明图记录与分布向量已经具有相同机制。

这样的构成定义还规定了错误处理的入口。缺少来源的候选只能作为假设记录，发生身份冲突时保留多个解释并请求新的观测，类型不一致时拒绝直接合并。我们不把非零向量视为理解的充分条件，而要求记录能参与任务读出、被证据修订并在更新中保持可追踪性。对只需要句子续写的系统，显式对象记录可能不是最经济的选择；对需要持续操作现实对象的系统，放弃身份和来源则可能使错误无法定位。语义单元由此成为表示契约，而不是从一种载体类比推出的普遍本体。

\section{形质绑定：关系兼容、候选竞争与提交条件}
\label{sec:section_5143}
\label{sec:section_8945}

结构分量允许多个关系同时成立，属性分量也允许多个字段同时存在。一个杯子可以同时位于桌面、靠近键盘并属于某个人；同一个杯子也可以同时是红色、陶瓷制成并具有较高温度。这种可组合性来自关系和属性的类型系统，不需要以玻色统计解释。另一方面，在规定了同一对象、同一时刻、同一分类准则的单选字段中，两个互斥类别不能同时作为已提交结果。互斥性来自任务契约，也不能从费米统计推出。多标签、复合对象、时间变化和分类体系重叠都允许不同的组合，因此我们必须先规定问题，再规定竞争方式。

我们把绑定理解为观测片段与对象记录之间的可检验归属。设当前有若干结构候选和若干属性观测，兼容矩阵的元素为零或一，分别表示类型、时间和来源约束是否允许该关联。对允许的关联，我们用无量纲评分比较位置连续性、内容匹配、观测可靠性及任务需求；不同分量先在各自尺度中处理，再通过明确权重组合。评分只是模型判断，不是观测事实。若所有候选均不兼容，系统应创建未知或冲突状态，而不是强制选出一个对象。遗漏这一分支会使归一化程序把本来无解的问题伪装成确定解释。

对第 $i$ 个结构候选，记非空合法集合为 $\mathcal J_i$，有限评分为 $s_{ij}$，无量纲探索尺度为 $T>0$。我们定义候选权重
\[
 p_{ij}=\frac{\exp(s_{ij}/T)}{\sum_{\ell\in\mathcal J_i}\exp(s_{i\ell}/T)},\quad j\in\mathcal J_i,
 \qquad p_{ij}=0\quad(j\notin\mathcal J_i).
\]
分母为正，故权重非负且和为一。对于合法候选 $k$，逐项使用商法则得到
\[
 \frac{\partial p_{ij}}{\partial s_{ik}}
 =T^{-1}p_{ij}(\delta_{jk}-p_{ik}).
\]
这说明提高一个候选评分会改变整行分配，而不是独立增加它的权重。较小尺度可使分配集中，但只要尺度正且评分有限，所有合法候选仍有正权重，归一化本身不会选出唯一赢家。权重是否可以解释为经过校准的概率，需要另行用数据检验；活动强度、评分和概率三者在接口中不能混用。

实际提交还需要单独的规则。对允许一对多归属的对象，我们可以保留多个属性观测，并记录各自来源；对要求一对一关联的跟踪任务，逐行最大值可能把两条轨迹都配给同一个观测，因而必须联合求解匹配，并允许未匹配状态。提交门槛可以同时要求评分差、来源可靠性及硬约束通过。门槛的数值属于工程设计，不能由归一化公式唯一推出。即使最佳候选被提交，也应保留备选和版本以便之后修订；提交是计算事务，不是语义实体发生了不可逆的自然变化。

红色圆形的例子揭示了绑定中最容易忽略的关系。如果一幅图同时存在红色桌面与圆形苹果，仅把红色向量和圆形向量合并，会形成一个没有证据支持的红苹果。我们要求颜色片段和形状片段共享可核验的对象关联，或者声明它们只是来自不同区域的组合假设。时间同步有助于限制候选，却不足以证明归属：两个不同对象可以在同一帧出现，同一对象的不同传感通道也可能存在延迟。若采用振荡编码，频率、相位和同步检测都必须另行定义；同步结果仍需要对象约束，不能凭相位一致推出知觉正确，更不能把绑定失败直接等同于临床诊断。

绑定的损失设计也需要尊重多个合法解释。一个观测区域可能同时包含杯子和杯柄，二者有整体—部件关系；若训练时把所有非唯一标签当成负例，模型会学会排斥真实的复合结构。我们可以分别训练身份关联、属性归属和关系预测，再检验联合结果是否满足一致性约束。联合损失下降说明所选指标改善，却不保证真实场景识别正确。遮挡、镜面反射、来源造假和训练分布之外的对象都可能使模型产生高分错误，故可靠绑定还需要异常检测和主动补充观测。

在运行中，我们把绑定建立、绑定保持和绑定撤销视为三种不同操作。建立时检查来源和类型，保持时检查时间连续性与新证据，撤销时把原结果标记为过期并保留审计历史。局部图提供可检查的对象位置，全局分布表示提供候选内容，两者通过带版本的关联表联合更新。任务从识别切换为操作时，绑定要求会更严格，因为错误归属可能导致错误执行。我们不要求所有智能都使用这一具体表结构，但任何声称保持对象连续性的实现，都需要说明错误绑定怎样被发现、怎样被纠正，以及纠正后哪些历史结论需要重新计算。

\begin{figure}[htbp]
    \centering
    \includegraphics[width=0.98\linewidth,height=0.48\textheight,keepaspectratio]{image/physound.png}
    \caption{原语义单元迁移图。图形主体保留，用于提示结构承载与活动传播的区别；其中声子、光子等历史名称仅为辅助意象，不构成绑定机制或统计规律的证明。}
\end{figure}

\section{迁移与近似不变量：从传播直觉到可核验约束}
\label{sec:semantion_duality_metaphor}
\label{sec:yuyizi_same}

我们需要区分三种经常被统称为传播的变化。第一种是身份相同的对象记录在不同表示坐标中重编码；第二种是活动沿固定关系结构扩散或被门控转移；第三种是新证据改变对象属性或类别。三者可以同时发生，却没有共同的无条件守恒量。把概念视为在介质中移动的意象有助于理解承载结构与活动的区别，但不能据此断言概念具有固定能量份额、思维速度服从某个色散公式，或语义流能够直接改变现实空间。实际硬件的通信延迟属于实现约束，应按测量值进入模型。

在表示迁移中，我们先规定需要保持的任务读出。设属性向量属于固定维数的实空间，两个坐标系统之间的可逆线性变换为 $R$，向量在新坐标中为 $q'=Rq$。若原读出是 $w^{\mathsf T}q$，则选择 $w'=R^{-\mathsf T}w$ 后，新读出为 $w'^{\mathsf T}q'=w^{\mathsf T}q$。相等来自两个变换互相抵消，不来自向量值不变。若原坐标采用正定度量 $G$，新度量应为 $G'=R^{-\mathsf T}GR^{-1}$，此时二次量同样保持。只变换向量而不变换读出与度量，一般不会保持任务意义。不可逆压缩则不能直接使用这个结论，需要另外给出读出误差界。

沿连续参数变化的坐标模型中，我们进一步规定 $q(t)$、正定度量 $G(t)$ 和矩阵 $A(t)$ 可微，时间按秒计，属性向量无量纲，$A$ 的单位为每秒。假设更新为 $\dot q=-Aq$，链式法则给出
\[
 \frac{d}{dt}(q^{\mathsf T}Gq)
 =q^{\mathsf T}\bigl(\dot G-A^{\mathsf T}G-GA\bigr)q.
\]
因而条件 $\dot G=A^{\mathsf T}G+GA$ 保证该二次量沿更新保持。这里我们从明确的微分条件得到数学结果，并未声称智能系统天然满足它。输入、遗忘或目标控制进入向量更新后会产生附加项；度量改变而不满足该条件时也不能继续宣称保持。对于离散实现，若采用 $q_{k+1}=P_kq_k$，精确保量需要 $P_k^{\mathsf T}G_{k+1}P_k=G_k$；一般的显式欧拉步并不自动满足这一式子。

活动传播是另一种模型。设固定连通无向图有非负对称权重，拉普拉斯矩阵为 $L$，活动向量为 $x(t)$，无外部输入时假设 $\dot x=-\kappa Lx$，其中 $\kappa$ 为每秒的更新率。由于 $\mathbf 1^{\mathsf T}L=0$，活动总和的导数为零；另一方面，$\frac d{dt}\|x\|^2=-2\kappa x^{\mathsf T}Lx\leq0$。前者保持总量，后者减少分布差异，不能互相替代。加入衰减或输入后总量一般变化；有向图也不能未经检查就保持算术总和。该模型说明一种结构约束下的传播机制，并不要求颜色、身份或逻辑真值像活动一样被平均。

我们把身份不变量限定为任务中的指称约束。例如，同一杯子沿桌面移动时，跟踪器希望保持身份键，却允许位置、姿态和温度变化。水果摊里的苹果与科技新闻里的苹果则可能是两个不同意义的对象，词形相同不构成核心特征保持的证据。国王和女王的关系变换可以保存某些角色约束，却可能改变性别字段、制度语境和指称对象，因此不能用向量旋转自动证明逻辑等价。有效迁移需要列出哪些字段必须保持、哪些允许改变、哪些需要新的观测，而不是寻找一个替所有任务定义意义的恒定模长。

类别稳定也不等于拓扑保护。在分类任务中，我们可以让扰动前后读出保持在同一标签区域，但这个区域只有在明确给定空间、边界和扰动集合后才可分析。身份变化、材料转化或分类体系修订都可能使类别合法改变。图的连通分量表示当前路径可达性，不等于一个概念永远不能跨越的范畴；路径同伦类需要连续空间与端点条件，也不能把它直接等同于同调或分类标签。我们允许模型使用拓扑工具，但要求每个不变量都有定义和保持条件，不能用同一个名称代替不同层次的性质。

迁移的验证因此采用分项测试。我们在固定变换集合下检查任务读出误差，在固定图模型下检查传播总量及数值步长，在对象轨迹上检查身份切换率，在属性更新上检查来源与约束。若指标改善，只能说明这些被测条件下的表现改善；不能推出全局分布表示与环境整体同构。某个内部量保持也不保证任务成功，例如错误绑定的身份键可能一直保持而实际跟踪了另一个对象。真正有用的连续性是可修订、可解释且与外部反馈相接的连续性，而不是内部数值始终不变。

\section{有限分辨率编码与语义单元的属性谱系}
\label{sec:section_1867}
\label{sec:section_5195}

连续输入具有远多于当前任务需要的细节。我们处理一段声音时不必保存传感器可分辨的所有采样误差，机器人抓取杯子时也不必在每次控制中重新表达整张图像。离散化首先是有限计算预算下的编码选择：系统保留足以支持任务的区别，合并暂时无关的变化。连续变量本身可以在数字设备上用有限精度近似，离散标签也可以伴随连续置信和未知状态。因此有限资源不推出所有概念天然不可分割，也不能推出理解必须一次激发一个完整概念。部分信息和逐步确认是正常的计算状态。

我们把编码器定义为从观测空间到有限码集合的映射，把重建器定义为从码集合返回任务表征的映射。设归一化观测空间为 $\mathcal Y\subset\mathbb R^m$，有限码集合为 $\mathcal C$，编码与重建分别为 $e$ 和 $d$。任务映射 $f:\mathcal Y\to\mathbb R^r$ 具有 Lipschitz 常数 $K$，即任务读出的差不超过观测差的 $K$ 倍。若有效观测范围内满足 $\|y-d(e(y))\|\leq\varepsilon$，则直接应用条件得到
\[
 \|f(y)-f(d(e(y)))\|\leq K\varepsilon.
\]
这是任务读出误差界，不是一般意义保存定理。归一化之前的长度、颜色或时间尺度应记录在编码契约中。若任务读出在边界附近不连续，例如阈值触发动作，这个界就不能未经调整继续使用；系统需要安全裕量、未知状态或更细分辨率。

对分类读出还可以明确边界条件。设两个类别评分函数的 Lipschitz 常数分别为 $K_1$ 和 $K_2$，原观测中第一类领先的差为 $m_0>0$。重建误差至多为 $\varepsilon$ 时，两评分差的变化幅度至多为 $(K_1+K_2)\varepsilon$。因此当 $m_0>(K_1+K_2)\varepsilon$ 时，这两类的次序保持；多类情形需要对所有竞争者检查。若领先差太小，我们不强行归类，而是提高采样精度或返回未决状态。此推导表明，离散标签的可靠性依赖观测误差和决策边界，而不是依赖一种语义能量阈值。

编码器的工程目标可以包含任务失真、存储量和处理延迟，但三个量必须各自计量。任务失真按规定的指标定义，存储按比特或字节定义，处理延迟按秒定义；实际功耗和累计焦耳能耗可以作为另一个预算分量，不能借用活动强度替代。若把这些分量组合为无量纲目标，需要显式参考尺度和权重，且不同权重代表不同设计偏好。优化所得码本不自动具有可解释的对象身份，所以下游仍需检查结构归属和来源。我们可以把 VTE 作为一种变分拓扑编码器的工程名称，但必须说明它输出什么、哪些约束如何训练，以及拓扑信息是否真的参与计算。

形质分离在采样流程中意味着先区分位置、时间、角色和属性，再声明它们如何联合，并不要求两种信息统计独立或正交。声音的时序结构与音色相关，图像的形状与纹理也可能共同指示对象，强制独立反而会丢失任务线索。一个可行流程是先生成候选片段和时间戳，再估计结构关系与属性，随后进行绑定、合并和版本提交。采样率、窗口长度及传感延迟决定哪些事件可被区分；结构编辑改变对象数量时，编码维度和关联表应一起更新。只在旧维度上修改向量可能留下指向已删除对象的悬空属性。

我们据此重建六类属性谱系。类别说明记录承担对象、关系、事件还是控制角色；定域说明它属于哪个结构位置、时间窗口和来源；内容说明可读出的属性及类型；活动说明当前参与计算的程度；修订代价说明修改、重新验证和传播影响所需资源；耦合说明在规定接口中对其他状态更新的影响。这六类不是六个普遍独立的物理量，既可以是标量，也可以是集合、向量或规则。尤其修订代价不能由度量行列式直接决定，耦合能力也不能从一个注意力权重推断出对象的全部作用。

以家庭成员记录为例，姓名和身份指称可以稳定，近期情绪和所在位置可以变化，证据可靠性需要随来源更新。修订一段错误描述的直接写入代价可能很低，但若该描述已经影响多项计划，重新检查后续结果的代价会很高。我们因此把存储修改成本、依赖传播成本和现实执行成本分开。冻结模型参数不使所有运行时记录不可修改，活动较高也不意味着记录更难修订。这种属性组织把原来的质量、能量和荷的意象转化为可测的工程字段，同时保留了它们所要解释的稳定性、参与程度和相互影响问题。

\section{跨系统比较与 AgentOS 中的工程落点}
\label{sec:section_2084}

跨系统比较应从任务接口出发，而不是预先给不同系统划分高低本体等级。我们关心一个系统如何区分对象与属性、如何建立关联、如何保持身份、如何修订以及如何使用外部反馈。不同实现可以在某些指标上相似，也可以在别的指标上差异明显。相同名称的活动量可能具有不同单位和含义，不能把神经发放率、模型注意力概率、环境信号浓度并列后直接比较大小。比较首先需要任务、时间尺度、数据来源和评价协议对齐，然后才讨论模型机制的可借鉴部分。

对生物认知，我们把视觉、听觉、触觉与行动反馈看作可用于建立和纠正绑定的多类信息来源。具体脑网络是否实现了某种张量积、哪类区域承担身份维持，以及这些机制是否可由本章模型解释，都属于需要经验材料检验的问题。一个作者提出的工程分解不能靠脑区名称获得生物学证明。我们可以设计对象错配、时间延迟、遮挡和跨模态冲突任务，比较模型的错误模式与观测，但相似错误模式仍然不足以确认唯一机制。这里保留生物系统作为约束来源，而不把它改写成我们架构的自然实现说明。

对语言模型，我们不把位置编码称为虚假的结构，也不把分布表示称为只有属性没有关系。文本位置、注意力、训练参数和工具接口都可能承载不同形式的关系信息，具体能力取决于模型及使用方式。更实际的问题是：某个部署是否提供持续对象身份、是否保留来源、是否具有跨回合状态，以及错误能否通过反馈修订。纯序列生成中，指称变化可能难以追踪；结合显式记录或检索接口后，这种限制可以被部分缓解。我们把幻觉分析为生成结果与可核验证据或约束不一致，不能用缺少一种指定形元解释所有幻觉，更不能据此否定其他实现的结构能力。

对群体系统，环境中留下的痕迹可以承担外部记忆，局部个体可以通过痕迹、位置和交互规则形成路径偏好。信号变化的单位、更新率与空间范围应按对象实际测量，而不是当作语义活动量的自然对应。某一群体在特定任务中缺少跨情境抽象接口，不能推出所有群体永远无法抽象；同样，形成有效路径也不说明群体已经拥有显式概念记录。我们可以比较任务迁移、环境变化后的恢复和痕迹失效的影响，但应明确系统级表现与个体内部机制之间的区别，避免用宏观结果反向指定唯一微观结构。

对理想智能的工程要求，我们把显式联合表示、动态修订和目标相关耦合看作可选择的设计原则，而不是实现通用智能的充分必要条件。需要长期操作对象的系统通常应能够追踪身份与证据，需要适应变化的系统应有某种有效更新通道，需要处理多目标冲突的系统应明确控制与边界。更新通道可以位于参数、记录、工具、技能或外部环境中，不要求全部表现为同一种在线参数学习。对象记录可以以图、表、向量索引或混合形式实现，只要接口可检验，便可纳入 HSF-HD 的一般状态动力学模型族。

在 AgentOS 实例中，我们把当前任务可用的绑定记录放入工作记忆 $h(t)$，把观测、提交、纠正和执行结果作为情景记忆 $E$ 的事件，把可跨情境复用的结构与更新规则整理入结构记忆 $\Theta$。外部记忆 $M_e$ 保存可检索的证据和较大的对象库，访问结果进入当前认知界面时携带来源与版本。自我相关状态 $\Psi$ 可以影响哪些身份与承诺需要优先保持，SPC 压缩运行状态，TCE 调节候选探索与增益，CCM 管理工作界面的资源负荷，Boundary 约束输入输出，Offloading 改变信息与技能的存放位置。这些部件共同提供工程接口，并不成为所有智能系统必须采用的唯一分解。

一次机器人桌面整理任务可以把这些接口串联起来。输入阶段产生杯子与杯柄的候选记录，绑定阶段核对它们的整体—部件关系，行动前检查抓取对象与温度测量是否指向同一身份。若遮挡造成身份不确定，控制器增加观测而不是把最大评分直接转成抓取动作；若外部数据库提供材料信息，记录同时保存其版本和来源；若执行后杯子位置发生变化，工作记忆更新位置，情景记忆保存动作结果，结构记忆只在重复证据足够时更新策略。这样，记录形成、内部推断、决策提交和现实执行有明确的边界，内部目标改善也不会被误报为任务完成。

验证方案应覆盖正确样例和反例。我们检查外观相同而身份不同的对象、类别相同而属性冲突的记录、单个来源被重复计数的情况、观测时延造成的错配，以及码本压缩丢失关键边界的情况。指标分别报告身份稳定、属性归属、校准、任务读出误差、修订延迟和执行后果；不能用一项总准确率代替全部机制审查。通过这些验证，语义单元的概念才从命名进入可检查的计算设计。系统论、协同论、耗散理论、复杂系统理论与认知科学提供一般原则和经验约束，HSF-HD组织状态与更新模型，AgentOS实现具体接口；三层关系在这一微观表示问题上仍然保持清楚。
